\begin{proposition}[Triangular freedom of reciprocal squares]
\label{prop:hasse-square-no-go}
Let $K$ be a field of characteristic different from~$2$, let $d\ge1$, and
put
\[
 \Delta(t)=t^2-34t+1.
\]
For a normalized reciprocal polynomial
\[
 B(t)=1+x_1t+\cdots+x_{d-1}t^{d-1}+t^d,
 \qquad x_{d-j}=x_j,
\]
and $\varepsilon\in\{0,1\}$, write
\[
 \Delta(t)^\varepsilon B(t)^2
 =\sum_{j=0}^{2d+2\varepsilon}H_jt^j.
\]
If $m<d/2$, then the coefficient map
\[
 (x_1,\ldots,x_m)\longmapsto(H_1,\ldots,H_m),
\]
with the remaining free coefficients of~$B$ held fixed arbitrarily, is a
triangular polynomial automorphism of $K^m$.  In particular, every vector
in $K^m$ occurs as $(H_1,\ldots,H_m)$ for some normalized reciprocal
polynomial~$B$.

Consequently, the factorization
\[
 H(t)=\Delta(t)^\varepsilon B(t)^2
\]
and reciprocity alone impose no restriction on a zero pattern contained in
the first $m<d/2$ coefficients of~$H$: any such pattern, together with its
reciprocal reflection, can occur.
\end{proposition}

\begin{proof}
For $1\le j\le m$, the inequality $j<d/2$ ensures that reciprocity has not
identified $x_j$ with any of $x_1,\ldots,x_{j-1}$.  The coefficient of
$t^j$ in the square is
\[
 [t^j]B(t)^2
 =2x_j+\sum_{i=1}^{j-1}x_ix_{j-i}.
\]
Since $\Delta(0)=1$, multiplication by $\Delta^\varepsilon$ gives
\[
 H_j=2x_j+Q_j(x_1,\ldots,x_{j-1})
\]
for a polynomial $Q_j\in K[x_1,\ldots,x_{j-1}]$; coefficients of~$B$
with index greater than~$j$ do not occur.  Thus the displayed map is
triangular with diagonal entries~$2$.  Given arbitrary
$(a_1,\ldots,a_m)\in K^m$, one obtains its unique preimage recursively from
\[
 x_j=\frac{a_j-Q_j(x_1,\ldots,x_{j-1})}{2}.
\]
After choosing any still-free middle coefficients, set $x_{d-j}=x_j$ to
complete~$B$.  This proves surjectivity and the polynomial inverse formula.

Both $B^2$ and~$\Delta^\varepsilon$ are reciprocal, so their product is
reciprocal of degree $2d+2\varepsilon$:
\[
 H_{2d+2\varepsilon-j}=H_j.
\]
Prescribing $a_j=0$ on an arbitrary subset of $\{1,\ldots,m\}$ therefore
realizes that zero pattern and its reflected copy.  The conclusion concerns
the full abstract family of reciprocal squares; it does not assert that the
arithmetically distinguished square root attached to the Ap\'ery motive is
generic in that family.
\end{proof}
